\section{System Architecture}
\label{sec:arch}
\system{} integrates multi-agent reasoning with industrial supervision around a common representation of production assets and operating context. Its organizing object is the node: an acquisition or control endpoint with engineering semantics, device connections, line membership, and agent bindings. A channel organizes the agents that interpret these endpoints, controllers mediate plant interactions, and Fig.~\ref{fig:arch} locates these components and execution boundaries.

\begin{figure*}[!t]
\centering
\includegraphics[width=0.9\textwidth]{figures/publication/fig2-architecture.pdf}
\caption{System architecture and execution boundaries. The co-resident platform core connects agent channels, semantic nodes, industrial controllers, and operator services. CLI/OMP engines and storage services extend beyond that process. Shared identifiers support attribution across the logical data model; they do not imply a distributed transaction or atomic physical actuation.}
\label{fig:arch}
\end{figure*}

\subsection{A Shared Supervisory Core}
The Node.js core hosts channel routing, task scheduling, industrial services, protocol adapters, and platform interfaces: REST and WebSocket endpoints serve operator applications, MCP and A2A interfaces support external integration, and the plugin host extends platform capabilities. These surfaces reuse platform objects, but their exposed operations and permission checks differ: an agent's industrial authority derives from its node bindings, not from channel access or from the mere availability of a protocol endpoint.

Co-residency reduces the need to synchronize duplicate nodes, bindings, and runs across separately assembled applications, but does not place the deployment in one process: CLI engines and the platform's bundled OMP agent harness can execute in spawned processes, TimescaleDB and MinIO are external, and connection establishment, queued protocol operations, and readback introduce asynchronous boundaries. Fast regulation and protection therefore remain with the plant controller and its safety equipment.

\subsection{Reasoning, Plant Interaction, and Operational Context}
On the reasoning side, a channel combines a lead, worker instances, a task tree, messages, and retrievable memory. The scheduler couples model-based supervision to deterministic dispatch and progress tracking, and engine adapters allow different reasoning implementations to participate without changing a bound node's identity or the controller interface. This separates task organization from device authority: assigning a worker a task does not grant permission to execute the associated industrial action.

On the plant side, acquisition nodes turn protocol readings into engineering quantities, control nodes translate engineering setpoints into device commands and readback results, per-node runtimes manage timing and in-flight operations, and controller services supply gateway state, recipe context, and write bookkeeping. Nodes and devices have many-to-many configuration associations: one device can supply several semantic nodes, and a node can reference several devices. Association is distinct from execution routing: a control-node write invokes its configured driver and connection settings rather than iterating over its associated devices, so the interface follows process meaning without implying multi-device actuation by one write.

The bridge between these sides consists of semantic cards, instance-specific bindings, industrial tools, and source-dependent controller checks; cards expose units, ranges, current values, and relevant process context. Bindings determine whether a channel member can request an operation and whether human approval is needed. Controllers then apply the checks appropriate to that request's origin: semantic guidance helps an agent propose a sensible action, whereas executable checks determine which requests reach a driver. Prompt guidance, such as the suggested step, is not an enforced interlock.

\subsection{Logical Continuity Across a Production Run}
Lines, products, recipes, and runs provide the context in which measurements and interventions become comparable; their identifiers connect platform state, time series, image frames, operational events, and optimization records across heterogeneous stores---SQLite and file-backed repositories for platform state, TimescaleDB for measurements, MinIO for objects. Together these constitute one logical data model, not one transaction domain: the journal records attributable changes and judgments, but is not a cryptographically signed ledger, and bookkeeping failure does not necessarily undo an already completed write.

An end-to-end interaction begins when an engineer connects devices, defines nodes, associates production context, and binds channel members; starting a line establishes an active run and initiates its parameter writes. Agents inspect observations, exchange evidence, and propose adjustments through their bound tools; accepted requests proceed to device communication and readback, and subsequent observations inform an agent-supplied judgment and an independent rule-based backstop; recovery may require a further physical write, not merely a recipe redefinition.

Operators participate through approvals, line operations, and monitoring. The device-twin renderer visualizes state without performing protocol I/O, but controls in the Town operator web interface can request API writes, so visualization is separated from device transport, not from all operator actuation. Section~\ref{sec:mech} details the six mechanisms and the cross-cutting HITL boundary, and marks where their guarantees end.
